\chapter{SYSTEM ANALYSIS}
 
 \section{Module Description}
 The "Bloomcycle" application is divided into multiple functional modules:
 \begin{itemize}
    \item \textbf{Authentication \& Authorization Module:} This module is responsible for secure user registration and login via JSON Web Tokens (JWT).
    \item \textbf{Period Tracking Module:} Core module logging start/end dates of menstrual cycles to compute history.
    \item \textbf{Symptom Logging Module:} Tracks daily physical/emotional symptoms (cramps, fatigue, mood) against specific dates.
    \item \textbf{Prediction \& Insights Module:} Predicts next cycle using historical averages and presents data through interactive dashboard graphs.
\end{itemize}

\section{Analysis of Architecture}
{\large \textbf{MERN Stack Architecture}}\\\\
Bloomcycle is built on MongoDB, Express.js, React.js, and Node.js.
\begin{enumerate}
    \item \textbf{React.js (Frontend):} Fast, component-based UIs displaying calendars and forms.
    \item \textbf{Node.js \& Express.js (Backend):} Fast API processing, JWT auth middleware.
    \item \textbf{MongoDB (Database):} Flexible NoSQL schema holding periods and variable symptom logs.
\end{enumerate}

\section{Feasibility Analysis}
\subsection{Technical Feasibility}
The system is highly technically feasible due to widespread support for MERN and cloud-based deployments.
\subsection{Economical Feasibility}
The economic feasibility is viable, utilizing minimal running costs via free tier cloud services like MongoDB Atlas.
\subsection{Operational Feasibility}
Operationally, Bloomcycle naturally integrates into standard daily tracking habits on web platforms.